\begin{abstract}
\sid{
Each paper must be no more than 8 pages (including the abstract, figures, and tables),double-columned, 9pt or 10pt font. One page of references is allowed, which does not count towards this 8-page limitation.
}\sid{(placeholder)} 

% Hardware/software (HW/SW) partitioning is an important problem in embedded system co-design, requiring tasks to be assigned to hardware or software under area constraints while minimizing makespan, i.e., the overall execution time from start to finish.
% HW/SW partitioning is NP-hard and tightly coupled with scheduling, since the execution order of tasks directly affects the final makespan.
% %
% Existing approaches, including mixed-integer linear programming (MILP), metaheuristics, and prior graph-based learning methods, either face scalability limitations, lack optimality guarantees, or do not jointly optimize partitioning and execution ordering in a unified differentiable framework, limiting their ability to consistently find high-quality solutions.
% %
% We present \diffgnn, a differentiable graph neural network (GNN) approach for joint HW/SW partitioning and scheduling. The proposed model combines a GNN encoder with a placement head for assignment prediction and an ordering head that learns execution priorities through a differentiable relaxation.
% Unlike prior approaches, \diffgnn embeds task-order learning into the optimization itself, enabling end-to-end learning of both placement and scheduling decisions in a unified manner. We introduce a soft scheduling surrogate that approximates makespan under hardware-area constraints.
% %
% We further incorporate a post-processing and refinement module that converts relaxed predictions into feasible discrete assignments and improves final solution quality. A practical strength of \diffgnn is that it is optimized directly on the given task graph, eliminating the need for a separate corpus of training graphs.
% %
% Across task graphs derived from neural network workloads, \diffgnn achieves lower makespan than metaheuristics and prior GNN-based HW/SW partitioning approaches, while remaining competitive with MILP on small and medium instances. In contrast to MILP, however, \diffgnn remains practical on larger graphs where exact optimization becomes computationally prohibitive. These results show that differentiable joint placement and order learning provide a scalable and effective new direction for HW/SW partitioning.
% Code is available at \url{https://github.com/anonymousauthors001/diff_gnn}
\end{abstract}

